\documentclass[12 pt, letterpaper]{hmcpset}
\usepackage{hw-preamble}

\name{Jayden Thadani}
\class{MATH 171 --- Abstract Algebra I}
\assignment{Exam 2}
\duedate{2nd April 2024}

\begin{document}

\begin{problem}[1.]
	Let $A$ and $B$ be subsets of $S_{4}$ where $A = \{ (1 \, 2), (3 \, 4) \}$ and $B = \{ (1 \, 2)(3 \, 4) \}$. Prove or disprove that $N_{S_{4}}(A)$ is a subgroup of $N_{S_{4}}(B)$.
\end{problem}

\begin{solution}
	\[
		\boxed{
			N_{S_{4}}(A) \le N_{S_{4}}(B).
		}
	\] \\
	Suppose $\sigma \in N_{S_{4}}(A)$. Recall that conjugation by any group element is an automorphism, and thus, the restriction of conjugation by $\sigma$ to just the set $A$, must still be bijective. \\\\
	That is, if $\sigma (1 \, 2) \sigma^{-1} = \alpha$ and $\sigma (3 \, 4) \sigma^{-1} = \beta$, we have $\alpha, \beta \in A$ and $\alpha \neq \beta$. Thus,
	\[
		\begin{split}
			\sigma (1 \, 2) (3 \, 4) \sigma^{-1} & = \sigma (1 \, 2) \sigma^{-1} \sigma (3 \, 4) \sigma^{-1} \\
							     & = \alpha \beta \\
							     & = \begin{cases}
								     (1 \, 2) (3 \, 4) \in B & \text{or} \\
								     (3 \, 4) (1 \, 2) \in B
							     \end{cases}
		\end{split}
	\]
	where $(3 \, 4)(1 \, 2) \in B$ because $(3 \, 4)(1 \, 2) = (1 \, 2)(3 \, 4)$, because disjoint cycles commute. \\\\
	Thus, any $\sigma$ that normalises $A$, normalises the only element in $B$. Thus, $N_{S_{4}}(A) \subset N_{S_{4}}(B)$. \\\\
	Since $N_{S_{4}}(A)$ is a group with the same operation as $N_{S_{4}}(B)$, it being a subset gives us $N_{S_{4}}(A) \le N_{S_{4}}(B)$.
\end{solution}

\pagebreak

\begin{problem}[2.]
	Prove that $S_{4}$ does not have a normal subgroup of order $8$.
\end{problem}

\begin{solution}
	Suppose $H \trianglelefteq S_{4}$ with $\lvert H \rvert = 8$. First we will show that $H$ cannot contain all of the transpositions in $S_{4}$, and next, we will show that if it does not, then a contradiction arises. \\\\
	Suppose $H$ contains all the transpositions in $S_{4}$. Note that there are $\binom{4}{2} = 6$ distinct transpositions in $S_{4}$. Then $(1 \, 2), (1 \, 3), (1 \, 4)$ are in $H$. Then $(1 \, 2) (1 \, 3) = (1 \, 3 \, 2)$, $(1 \, 3) (1 \, 2) = (1 \, 2 \, 3)$, and $(1 \, 4) (1 \, 2) = (1 \, 2 \, 4)$ are all distinct elements of $H$. We can see that they are distinct since $2$ is mapped to a different number by each of them, and they must be in $H$ since $H$ is closed under group multiplication. \\\\
	This means that we already have $9$ distinct elements that must be in $H$, contradicting the fact that $H$ has only $8$ elements. Thus, our assumption that $H$ contains all transpositions in $S_{4}$ must be incorrect. There must exist some transposition in $S_{4}$ that is not in $H$. \\\\
	Suppose $\sigma$ is one such transposition in $S_{4}$ with $\sigma \notin H$. Then $K = \left < \sigma \right >$ is a subgroup of order $2$ in $S_{4}$, with $\lvert H \cap K \rvert = 1$ since $H$ does not contain the only non-trivial element of $K$. Note that since $H$ is normal, $HK = KH$ and thus, $HK$ must be a group. However, by the formula for the order of a product,
	\[
		\begin{split}
			\lvert HK \rvert & = \frac{\lvert H \rvert \lvert K \rvert}{\lvert H \cap K \rvert} \\
					 & = \frac{8 \times 2}{1} \\
					 & = 16.
		\end{split}
	\]
	Since $16$ does not divide $\lvert S_{4} \rvert = 24$, $HK$ being a subgroup contradicts Lagrange's theorem, and thus, our assumption that there exists a normal subgroup of order $8$ is incorrect.
\end{solution}

\pagebreak

\begin{problem}[3.]
	Write up a proof of the 3rd isomorphism theorem.
\end{problem}
	\begin{theorem}
		For any normal subgroups $H, K$ of $G$, with $H \le K$, we have
		\[
			K/H \trianglelefteq G/H
		\]
		and
		\[
			(G/H)/(K/H) \cong G/H. 
		\]
	\end{theorem}
	\begin{proof}
		We claim that the map
		\[
			\begin{split}
				\varphi \colon & G/H \to G/K \\
					       & gH \mapsto gK
			\end{split}
		\]
		is a surjective homomorphism with kernel $K/H$ (and that thus, by the first isomorphism theorem we have $K/H \trianglelefteq G/H$ and $(G/H)/(K/H) \cong G/K$). \\\\
		First we show that $\varphi$ is well-defined. Suppose $g_{1} H = g_{2} H$. Then $g_{1} g_{2}^{-1} \in H$. Since $H \le K$, this gives us that $g_{1} g_{2}^{-1} \in K$, and thus $g_{1} K = g_{2} K$. That is, $\varphi(g_{1} H) = \varphi(g_{2} H)$, and thus, $\varphi$ is well-defined. \\\\
		Now we show that $\varphi$ is a homomorphism. To do so we can just compute $\varphi((g_{1}H)(g_{2} H))$ for any $g_{1}H, g_{2}H \in G/H$.
		\[
			\begin{split}
				\varphi((g_{1}H)(g_{2}H)) & = \varphi((g_{1}g_{2})H) \\
							  & = (g_1 g_{2}) K \\
							  & = (g_{1} K) (g_{2} K) \\
							  & = \varphi(g_{1} H) \varphi(g_{2} H).
			\end{split}
		\]
		This proves that $\varphi$ is a homomorphism. \\\\
		For any $gK \in G/K$, we know that there exists $gH \in G/H$, and by definition, $\varphi(gH) = gK$. Thus, $\varphi$ contains every $gK \in G/K$ in its image. That is, $\varphi(G) = G/K$. \\\\
		Finally, we will show that $\ker \varphi = K/H$ by double containment. \\\\
		Suppose $gH \in \ker \varphi$. Then $\varphi(gH) = 1_{G/K} = K$. But we know $\varphi(gH) = gK$. Thus, we must have $gK = K$, and thus $g \in K$. Then $gH \in K/H$. Thus $\ker \varphi \subset K/H$.  \\\\
		Suppose $kH \in K/H$ where $k \in K$. Then
		\[
			\begin{split}
				\varphi(kH) & = kK \\
					    & = kK \\
					    & = K \\
					    & = 1_{K/H}.
			\end{split}
		\]
		Thus, $kH \in \ker \varphi$. Thus, $K/H \subset \ker \varphi$. Then, by double containment. $K/H = \ker \varphi$. \\\\
		Finally, by the first isomorphism theorem $K/H \trianglelefteq G/H$ and $(G/H)/(K/H) \cong G/K$.
	\end{proof}

\pagebreak

\begin{problem}[4.]
	Let $H$ be a normal subgroup of a finite group $G$. Prove that if $G/H$ has an element of order $n$, then $G$ has an element of order $n$.
\end{problem}

\begin{solution}
	Let $\overline{g} = gH \in G/H$ be an element of order $n$. Thus, $\overline{g}^{j} \neq \overline{g}^{k}$ for any $j \neq k$ with $j, k \in \{ 1, \dots, n \}$ and $\overline{g}^{n} = H$. Thus, $g^{j} g^{-k} \notin H$ for all those $j, k$ and $g^{n} = h \in H$. \\\\
	Since $G$ is finite, $h$ must have finite order --- say $\lvert h  \rvert = m$. Then we have that
	\[
		\begin{split}
			(g^{m})^{n} & = (g^{n})^{m} \\
				    & = h^{m} \\
				    & = 1_{G}.
		\end{split}
	\]
	We claim that $g^{m}$ has order $n$. Suppose $(g^{m})^{l} = 1_{G}$ for some positive integer $l$. Then $g^{ml} = 1_{G}$. By the division algorithm, we have that $ml = nq + r$ where $0 \le r < n$. Then we can rewrite
	\[
		\begin{split}
			g^{ml} & = g^{nq + r} \\
			       & = (g^{n})^{q} g^{r} \\
			       & = h^{q} g^{r}.
		\end{split} 
	\]
	Since this is $g^{ml} = 1_{G}$, we have $h^{q} g^{r} = 1_{G}$, and thus, $g^{r} = h^{-q} \in H$. Since $n$ is the least positive number such that $g^{n} \in H$, we must have $r = 0$. Thus, $ml = nq$, and $h^{q} = g^{ml} = 1_{G}$. \\\\
	Since the order of an element divides any power equal to $1_{G}$, we have $m \mid q$, and thus, $q/m = d$ is an integer. Then $l = nd$ and thus, any positive $l$ such that $(g^{m})^{l} = 1_{G}$ is greater than or equal to $n$. In other words, $n$ is the least power such that $g^{m}$ raised to it is the identity, or $\lvert g^{m} \rvert = n$. \\\\
	Thus, $g^{m}$ is the element of $G$ with order $n$.
\end{solution}

\pagebreak

\begin{problem}[5.]
	Briefly explain why the following statements are true or false
	\begin{enumerate}
		\item If $K \trianglelefteq H$ and $H \trianglelefteq G$ then $K \trianglelefteq G$.
		\item If $H \le G$, then the left cosets of $H$ partition $G$.
		\item If $H \trianglelefteq G$, then $H$ is the kernel of some homomorphism with domain $G$.
		\item If $\lvert G \rvert < \infty$ and $n$ divides $\lvert G \rvert$, then $G$ has a subgroup $H$ such that $\lvert H \rvert = n$.
		\item If $H$ and $K$ are subgroups of $G$ then $HK$ is a subgroup of $G$.
	\end{enumerate}
\end{problem}

\begin{solution}
	\begin{enumerate}
		\item \fbox{False} \\\\
			From the conditions we have that $K$ is invariant under conjugation by any $h \in H$ and $H$ is invariant under conjugation by any $g \in G$. However, it is still possible that we could have some $g \in G$ such that $gkg^{-1} = h \in H \setminus K$ for some $k \in K$. Then $gkg^{-1} \notin K$ and thus, $K$ would not be normal. \\\\
			From Dummit \& Foote (page 91) we have the specific counterexample that $\left < s \right > \trianglelefteq \left < s, r^{2} \right > \trianglelefteq D_{8}$ since each subgroup has index $2$ in the next. ($\lvert \left < s \right > \rvert = 2$, $\lvert \left < s, r^{2} \right > \rvert = 4$ and $\lvert D_{8} \rvert = 8$.) However, $\left < s \right >$ is not a normal subgroup of $D_{8}$ since $rsr^{-1} = sr^{2}$ which is not in $\left < s \right >$. Notice that, as we described, $rsr^{-1}$ is still a member of $\left < s, r^{2} \right >$, but is not one of $\left < s \right >$.
		\item \fbox{True} \\\\
			We will prove that the relation $g_{1} \sim g_2 \iff g_{1} \in g_{2}H$ is an equivalence relation, and thus, the left cosets of $H$ are equivalence classes of $G$, which must partition the set. Let $g_{1}, g_{2}, g_3 \in G$. \\\\
			We know that $g_{1} \in g_{1} H$ since $1_{G} \in H$ and thus $g \, 1_{G} = g_1 \in g_{1}H$. \\\\
			If $g_{1} \sim g_2$, then $g_{1} \in g_{2} H$ giving us $g_{1} = g_{2} h$ for some $h \in H$. Right-multiplying by $h^{-1}$ on both sides gives us $g_{1} h^{-1} = g_{2}$ and thus $g_{2} \in g_{1} H$. Thus $g_{2} \sim g_{1}$. \\\\
			Suppose $g_{1} \sim g_{2}$ and $g_{2} \sim g_{3}$. Then $g_{1} = g_{2} h_{1}$ and $g_{2} = g_{3} h_{2}$ for some $h_{1}, h_{2} \in H$. Thus $g_{1} = g_{3} h_{2} h_{1}$ and $g_{1} \sim g_{3}$. \\\\
			Thus, $\sim$ is reflexive, symmetric, and transitive --- an equivalence relation. The equivalence class of $g_{1}$ under this equivalence relation is by definition $g_{1}H$ since $g_{2} \sim g_{1}$ only if $g_{2} \in g_{1} H$. Since the left cosets of $H$ are all equivalence classes, they partition the group $G$.
		\item \fbox{True} \\\\
			Since $H \trianglelefteq G$, $G/H$ is a group with the standard binary operation. Consider the projection $\pi : G \to G/H$ by $g \mapsto gH$. Notice that $gH = H = 1_{G/H}$ if and only if $g \in H$. Thus $\ker \pi = H$.
		\item \fbox{False} \\\\
			We have shown in class that the group of rotational symmetries of the regular tetrahedron has order $12$, but no subgroup of order $6$, even though $6 \mid 12$. A proof is provided in Dummit \& Foote (on page 92). This is one example of how the converse of Lagrange's theorem is false.
		\item \fbox{False} \\\\
			We have shown in class (in the lemma prior to our proof of the second isomorphism theorem) that $HK$ is a subgroup of $G$ if and only if $HK = KH$. We will show that there exists a group $G$ with subgroups $H$ and $K$ such that $HK \neq KH$. \\\\
			If we consider the free group $G = \left < h, k \right >$ then we have subgroups $H = \left < h \right >$ and $K = \left < k \right >$. Then $HK \neq KH$ since any $g_{1} \in HK$ is of the form $g_{1} = h^{m} k^{n}$ for some integers $m, n$ which by the definition of the free group is not equal to $g_{2} = k^{l} h^{j}$ for any integers $l, j$, which is the form of any element of $KH$. Thus, $HK$ is not a group. We can also see this in that $(hk)^{-1} = k^{-1}h^{-1}$ is not in $HK$ even though $hk$ is. \\\\
			Note to instructor: I know we haven't actually covered free groups and from reading Dummit \& Foote it seems like they need a lot of machinery to define, but I hope this is okay as a counterexample.
	\end{enumerate}
\end{solution}

\end{document}
